% Irodov Part VI §6.3, original 97--104, printed page 245.
% Input: user-provided PDF page 254 and appendix Rydberg constant.
\atomictermsproblem{97}{245}
锂原子的价电子在$2S$及$2P$态的结合能
分别为$5.39\,\mathrm{eV}$和$3.54\,\mathrm{eV}$。
求对应里德伯修正量。
\exsolution
碱金属结合能近似为
$E_{nL}=\mathrm{Ry}/(n+\alpha_L)^2$，
其中$\mathrm{Ry}=13.606\,\mathrm{eV}$。
本题$n=2$，故
\[
 \boxed{\alpha_S=\sqrt{13.606/5.39}-2\simeq-0.41},
 \qquad
 \boxed{\alpha_P=\sqrt{13.606/3.54}-2\simeq-0.04}.
\]
\excheck 采用$n+\alpha$记号故修正量为负，
常见量子亏损$\delta=-\alpha$为正。

\atomictermsproblem{98}{245}
钠价电子基态$3S$结合能为$5.14\,\mathrm{eV}$，
第一次激发电势为$2.10\,\mathrm V$。
求$3P$项的里德伯修正量。
\exsolution
激发态价电子结合能
$E_b(3P)=5.14-2.10=3.04\,\mathrm{eV}$，
故
\[
 \boxed{\alpha_{3P}=\sqrt{\mathrm{Ry}/3.04}-3
 \simeq-0.88}.
\]
\excheck 由于电子被激发后更容易电离，
$3P$结合能低于$3S$结合能。

\atomictermsproblem{99}{245}
锂锐线系第一条谱线波长$\lambda_1=813\,\mathrm{nm}$，
系列短波限$\lambda_2=350\,\mathrm{nm}$。
求基态价电子结合能。
\exsolution
锐线系对应$nS\to2P$。
短波极限给$E_{b,2P}=hc/\lambda_2$，
首条谱线$3S\to2P$给
$E_{b,3S}=hc(1/\lambda_2-1/\lambda_1)$。
S系列修正量
$\alpha_S=\sqrt{\mathrm{Ry}/E_{b,3S}}-3$，
所以
\[
 \boxed{E_{b,2S}=\frac{\mathrm{Ry}}{(2+\alpha_S)^2}
 \simeq5.3\,\mathrm{eV}}.
\]
\excheck 本题最终求的是$2S$基态结合能，
而非由系列短波限直接得出的$2P$结合能。

\atomictermsproblem{100}{245}
锂的S、P项里德伯修正量分别为$-0.41$、$-0.04$。
求由$3S$经$2P$至基态$2S$时发出的谱线波长。
\exsolution
电偶极跃迁要求$\Delta L=\pm1$，
故顺次发生$3S\to2P$、$2P\to2S$。
令$E_{b,nL}=\mathrm{Ry}/(n+\alpha_L)^2$，
\[
 \lambda_1=\frac{hc}{E_{b,2P}-E_{b,3S}},
 \qquad
 \lambda_2=\frac{hc}{E_{b,2S}-E_{b,2P}}.
\]
代入各$n,L$项，得
\[
 \boxed{\lambda_{3S\to2P}\simeq0.82\,\mu\mathrm m},
 \qquad
 \boxed{\lambda_{2P\to2S}\simeq0.68\,\mu\mathrm m}.
\]
\excheck 两光子能量的和等于初末能级差；
不能直接认为$3S\to2S$是一条允许电偶极谱线。

\atomictermsproblem{101}{245}
钠共振黄线双线的波长为
$589.00$及$589.56\,\mathrm{nm}$，
均对应$3P\to3S$，求$3P$项精细劈裂能。
\exsolution
因两线共用$3S$下能级，其光子能量之差
等于上能级的差：
\[
 \boxed{\Delta E=hc\left|
 \frac1{589.00\,\mathrm{nm}}-\frac1{589.56\,\mathrm{nm}}
 \right|\simeq2.0\times10^{-3}\,\mathrm{eV}}.
\]
\excheck 必须用倒波长差，而非波长的直接差。

\atomictermsproblem{102}{245}
铯锐线系第一线是双线，
波长为$1358.8$和$1469.5\,\mathrm{nm}$。
求后续谱线双线分量的角频率间隔。
\exsolution
锐线系同一末态$P$项存在恒定精细结构分裂，
因此全部锐线系双线角频率间隔一致。
\[
 \boxed{\Delta\omega=
 2\pi c\left|\frac1{1358.8\,\mathrm{nm}}
 -\frac1{1469.5\,\mathrm{nm}}\right|
 \simeq1.05\times10^{14}\,\mathrm{rad/s}}.
\]
\excheck 角频率比普通频率相差$2\pi$。

\atomictermsproblem{103}{245}
列出氢原子电子主量子数$n=3$时
所有可能的光谱项。
\exsolution
轨道量子数$l=0,1,2$，
电子自旋$s=1/2$，
耦合允许$j=l\pm1/2$，舍去负$j$。
于是
\[
 \boxed{3S_{1/2},\quad3P_{1/2},\quad3P_{3/2},\quad
 3D_{3/2},\quad3D_{5/2}}.
\]
\excheck $l=0$时只允许$j=1/2$，共五个谱项。

\atomictermsproblem{104}{245}
原子轨道角动量与自旋量子数$(L,S)$分别为
(a)$(3,2)$；(b)$(3,3)$；(c)$(2,5/2)$，
求总角动量量子数$J$的所有可能取值。
\exsolution
根据三角耦合规则，
$J=|L-S|,|L-S|+1,\ldots,L+S$，
\[
 \boxed{\text{(a)}\ J=1,2,3,4,5},
 \qquad
 \boxed{\text{(b)}\ J=0,1,2,3,4,5,6},
\]
\[
 \boxed{\text{(c)}\ J=\tfrac12,\tfrac32,\tfrac52,\tfrac72,\tfrac92}.
\]
\excheck 整数与半整数耦合时总$J$为半整数。
